\documentclass[11pt]{article}

\usepackage[a4paper,margin=0.78in]{geometry}
\usepackage{graphicx}
\usepackage{booktabs}
\usepackage{array}
\usepackage{float}
\usepackage[hidelinks]{hyperref}
\usepackage{xcolor}
\usepackage{caption}
\usepackage{subcaption}
\usepackage{enumitem}
\usepackage{amsmath}

\graphicspath{{figures/}}
\setlength{\parindent}{0pt}
\setlength{\parskip}{0.55em}
\setlist[itemize]{topsep=0.25em,itemsep=0.15em,leftmargin=1.3em}

\title{CW305 Power Side-Channel Input Recovery\\Formal Hardware Report}
\author{Khalifa University lab work notes}
\date{30 August 2026}

\begin{document}
\maketitle

\begin{abstract}
This report explains the current hardware result for my FPGA power
side-channel experiment. The goal is to reproduce the main idea of
Power2Picture and the active power-template attack, but on my own hardware
setup. I use the CW305 FPGA board and ChipWhisperer-Lite. The important result
is that the new data is not simulation and not only a host-PC reconstruction.
The FPGA is programmed with a custom leakage gateware, the output datapath is
checked against a CPU golden model, and the traces are captured from the board
using ChipWhisperer. From these traces I can recover recognizable MNIST-like
digit images. The best fresh active-template run reached 97.5\% recovered digit
recognition on 40 held-out images. The best Power2Picture-style generator
reached 90.33\% binary recognition and 93.0\% gray-output recognition on 300
test traces. I also ran a harder single-shot trained-kernel test. It reached
91.15\% pixel accuracy but only 45\% recognition, which shows the limit of the
low-cost measurement setup when averaging is removed.
\end{abstract}

\tableofcontents
\newpage

\section{INTRODUCTION}

The main question in this work is simple: can the power trace from an FPGA
neural-network accelerator leak enough information to reconstruct the input
image?

The reference papers show this idea using stronger measurement equipment and a
different FPGA setup. My bench is different. I use:

\begin{itemize}
  \item CW305 FPGA board,
  \item ChipWhisperer-Lite for the analog power capture,
  \item a custom binary 3 by 3 convolution gateware,
  \item MNIST digit images as the input data.
\end{itemize}

Because the hardware is different, I did not try to copy the paper line by
line. I kept the same research idea, but I adapted the experiment so it can run
on this board. The important rule in this report is that I only report recovery
numbers when the trace folder passes the hardware checks. The checks prove that
the FPGA was programmed with the custom design, that the board output matched
the software golden model, and that the traces came from the ChipWhisperer
capture path.

\begin{figure}[H]
  \centering
  \includegraphics[width=\linewidth]{fig01_status_summary.png}
  \caption{Current status after the reset and rerun. Normal result graphs use
  color. Only the flow-chart style diagrams are kept in black, white, and gray.}
  \label{fig:status}
\end{figure}

\section{REFERENCE PAPER}

The Power2Picture paper studies input recovery from power traces of neural
network accelerators on FPGA. The main idea is that early convolution layers
still carry strong information about the input image. If the power measurement
has enough detail, a model can learn the mapping from side-channel traces to
the original image.

The active power-template idea is related but uses a different route. It builds
a profile from known images and known power traces. During the attack, it uses
only the unknown trace and the profile to infer the local image patches. Then
overlapping patches are voted together to rebuild the full image.

In this work I use both ideas:

\begin{itemize}
  \item an active-template method, because it makes it easy to check if the
  trace alone carries local patch information;
  \item a Power2Picture-style generator, because it is closer to the learned
  trace-to-image approach.
\end{itemize}

The comparison is not one-to-one. The reference papers had better analog
conditions and different FPGA hardware. My implementation uses a custom leakage
core to raise the signal level enough for CW-Lite.

\begin{figure}[H]
  \centering
  \includegraphics[width=\linewidth]{fig07_reference_vs_our_setup.png}
  \caption{Main difference between the reference papers and this CW305
  implementation. The research question is the same, but the hardware and
  measurement quality are not the same.}
  \label{fig:reference-compare}
\end{figure}

\section{WHY THE FIRST IMPLEMENTATION WAS WEAK}

The first implementation was weak for five reasons.

First, the board identity problem was real. After programming, register pages
2, 3, and 4 were reading back as \texttt{0x02}, \texttt{0x05}, and
\texttt{0x2e}. These values match the stock CW305 AES design. In that state the
host PC can still run Python scripts, but the FPGA is not running the target
neural-network leakage design. Any recovery number from that state is unsafe to
claim.

Second, the first scripts did not stop hard enough when this happened. The new
scripts now reject stock AES before capture. This is important because a
software pipeline can make nice-looking files even when the FPGA is doing the
wrong job.

Third, the passive background detector was too simple for this gateware. It
assumed that background pixels always create lower activity. In the custom
binary convolution design, the visible leakage is closer to a match count
between a 3 by 3 patch and a kernel. This does not always make the background
the lowest-power case.

Fourth, the original paper had stronger power measurement conditions. CW-Lite
with the CW305 gives a useful signal, but it has less analog detail per clock
cycle. Clock activity, board noise, USB noise, and static current hide part of
the useful data.

Fifth, I found a split bug during the rerun. In the active-template script,
\texttt{--start-index} means the MNIST image index, not the offset inside the
captured folder. One precheck accidentally overlapped profile and attack
images. I removed that result and reran with non-overlapping image ranges.

\section{HARDWARE GATEWARE}

The FPGA design used for the final run is the custom leakage gateware
\texttt{cw305\_leakage\_top}. It is not the stock AES gateware. The bitstream
used in the fresh runs is:

\begin{itemize}
  \item file: \path{build/cw305_leakage_d8_l63.bit},
  \item SHA-256:
  \texttt{a5bb029e693174a3a6e63384ba5baa1f464de675c9dc1c16daf290087babf353}.
\end{itemize}

The design implements a binary 3 by 3 convolution over a 28 by 28 MNIST image.
It produces 26 by 26 valid output positions, so each image gives 676 window
outputs. The host writes the image and kernel into FPGA registers, starts the
core, waits for idle, and reads back the 676 signed scores.

The gateware also keeps a small leakage bank active during the computation.
This is an important difference from a normal accelerator. The purpose is not
to make a secure accelerator. The purpose is to make a controlled side-channel
victim that is measurable using the CW305 and CW-Lite setup. I would describe
this clearly in the paper, because it is an adapted experiment and not a direct
clone of the reference hardware.

\section{HARDWARE VALIDATION}

After the switch check and USB reset, I reran the board checks. The mode
switches were set as M0=1, M1=1, and M2=1. The custom bitstream was programmed
again. Then the identity readback gave:

\begin{quote}
\texttt{design\_signature\_pages\_2\_3\_4=[0, 0, 0]}\\
\texttt{PASS: 676 outputs bit-exact}
\end{quote}

This means the board was not answering like stock AES, and the neural-network
datapath matched the CPU golden model. This is the main protection against
accidentally training on traces from the wrong FPGA image.

I also ran one extra live challenge-response check after the harder capture,
without reprogramming the FPGA. I used MNIST image \texttt{1234} and trained
kernel \texttt{4}. The page readback was \texttt{[1, 0, 0]}, which is still not
the AES signature, and the check again printed \texttt{PASS: 676 outputs
bit-exact}. This matters because it checks a different image/kernel pair on the
already-programmed board. The page values are not a clean build-ID register in
this gateware, so I do not treat \texttt{[0, 0, 0]} as a fixed ID. I treat
``not the AES signature plus 676 correct outputs'' as the hardware proof.

\begin{figure}[H]
  \centering
  \includegraphics[width=0.88\linewidth]{fig03_register_identity_gate.png}
  \caption{Hardware identity check after reset. The readback is not the stock
  AES signature, and the 676-output functional check passes.}
  \label{fig:identity-check}
\end{figure}

The validation files also record the capture source. The fresh trace folders
all report:

\begin{itemize}
  \item \texttt{capture\_mode = live\_cw305\_chipwhisperer},
  \item \texttt{trace\_source = cw\_lite\_analog\_sync\_x4},
  \item ChipWhisperer version \texttt{6.0.0},
  \item scope serial \texttt{50203220415447303030313238323035},
  \item target serial \texttt{5020322030384a583330343234303030},
  \item ADC locked at the beginning and end of capture,
  \item functional check passed with zero mismatches.
\end{itemize}

This does not mean that the PC is not used. The PC is used for storing traces,
training, and reconstruction. But the measured signal itself is from the FPGA
power capture through ChipWhisperer. That is the same side-channel experiment
style as the reference paper: hardware produces the leakage, and the host does
the analysis after capture.

\section{CAPTURE METHOD}

The capture method uses the CW-Lite clock and ADC in a synchronized mode:

\begin{itemize}
  \item FPGA clock: 5 MHz from CW-Lite,
  \item ADC source: \texttt{clkgen\_x4},
  \item samples per FPGA cycle: 4,
  \item dwell per convolution window: 8 cycles,
  \item samples per window: 32,
  \item useful samples per image/kernel run: 676 $\times$ 32 = 21,632,
  \item total samples per trace: 22,144 with margin,
  \item trigger: FPGA trigger line to CW-Lite,
  \item gain: 40 dB.
\end{itemize}

For the active-template captures, I captured 9 kernel traces for every image.
In the easier run, each kernel trace was averaged over 5 repeated captures. In
the harder run, I used only 1 capture per kernel. For the Power2Picture-style
captures, I captured one single-shot trace per image using one trained kernel.

\begin{figure}[H]
  \centering
  \includegraphics[width=\linewidth]{fig16_rerun_trace_evidence.png}
  \caption{Examples of real captured power signal data. The upper part shows
  active-template segmentation and a power feature map. The lower part shows a
  P2P trace from a trace shard.}
  \label{fig:trace-evidence}
\end{figure}

\begin{figure}[H]
  \centering
  \includegraphics[width=\linewidth]{fig04_trace_to_features.png}
  \caption{How the raw trace is converted to per-window power features. This
  step is important because reconstruction works on features extracted from the
  captured side-channel trace.}
  \label{fig:trace-features}
\end{figure}

\section{METHODS USED AND IMPLEMENTATION}

\subsection*{Active power-template method}

The active-template method has two sides. The profile side uses known images
and their traces. It builds a table from a 3 by 3 image patch to a power feature
vector. The attack side receives only the trace-only bundle and the profile.
It does not receive the attacked image or the label.

For the clean one-hot run:

\begin{itemize}
  \item profile images: \texttt{0700--0739},
  \item attack images: \texttt{0740--0779},
  \item attack set size: 40 held-out images,
  \item kernel probes: 9 one-hot kernels,
  \item repeats: 5 captures per kernel.
\end{itemize}

For the trained-kernel run:

\begin{itemize}
  \item profile images: \texttt{0780--0819},
  \item attack images: \texttt{0820--0859},
  \item attack set size: 40 held-out images,
  \item kernel probes: 9 trained convolution kernels,
  \item repeats: 5 captures per kernel.
\end{itemize}

\begin{figure}[H]
  \centering
  \includegraphics[width=\linewidth]{fig05_algorithm_flow.png}
  \caption{Active-template attack flow. This flow chart is black and white on
  purpose. The truth image is outside the attack path and is used only after
  reconstruction for scoring.}
  \label{fig:active-flow}
\end{figure}

\subsection*{Power2Picture-style generator}

The Power2Picture-style method uses supervised learning. During training, it
uses pairs of known traces and known images. During testing, the model input is
only the trace feature vector. The truth image and label are loaded only after
the output exists, for scoring.

I tested three small variants on the same 2,000 hardware traces:

\begin{itemize}
  \item clock-level absolute features with gradient-MSE loss,
  \item window-level absolute features with gradient-MSE loss,
  \item clock-level absolute features with plain MSE loss.
\end{itemize}

The split was 1,500 training traces, 200 validation traces, and 300 test traces.
The generator script now has \texttt{--require-hardware}. If the manifest looks
like simulation or does not have hardware provenance, the run stops.

\section{HARDER HARDWARE TESTS}

To make the validation harder, I ran a new live capture after the previous
report. This test used trained kernels but removed repeat averaging:

\begin{itemize}
  \item folder: \path{host/traces_rerun_20260830_trained_avg1_120},
  \item images: 120 MNIST images,
  \item image range: \texttt{0900--1019},
  \item kernels: 9 trained convolution kernels,
  \item averaging: 1 capture per kernel,
  \item acquisition order: shuffled using seed \texttt{8303030},
  \item validation: passed,
  \item functional check: 676 outputs, zero mismatches.
\end{itemize}

This is harder because the previous active-template runs averaged 5 captures
per kernel. Removing this averaging makes the trace less stable. It also makes
the search problem harder, because the candidate sets become larger. I first
started a 60-image held-out attack, but the Python candidate search was taking
too long. I stopped that run and used a 20-image held-out subset so the report
contains a complete result. I do not report the interrupted 60-image attack as
a result.

\begin{figure}[H]
  \centering
  \includegraphics[width=\linewidth]{fig20_hard_avg1_trace_noise.png}
  \caption{The harder test removes repeat averaging. The trace still comes from
  the same ChipWhisperer/CW305 path, but the reconstruction becomes less
  stable.}
  \label{fig:hard-trace}
\end{figure}

\begin{figure}[H]
  \centering
  \includegraphics[width=0.62\linewidth]{fig21_rerun_hard_avg1_examples.png}
  \caption{Hard active-template examples with trained kernels and only one
  capture per kernel. The result is weaker but still shows some stroke
  recovery.}
  \label{fig:hard-examples}
\end{figure}

\section{RESULTS}

Table~\ref{tab:capture-sets} lists the new hardware trace sets. These are the
sets I would treat as the main evidence because they pass the hardware checks.

\begin{table}[H]
\centering
\caption{Fresh hardware trace sets used in this report.}
\label{tab:capture-sets}
\small
\begin{tabular}{>{\raggedright\arraybackslash}p{0.29\linewidth}cccc>{\raggedright\arraybackslash}p{0.22\linewidth}}
\toprule
Trace set & Images & Kernels & Avg. & Validation & Use \\
\midrule
\path{traces_rerun_20260830_onehot_80} & 80 & 9 one-hot & 5 & pass & active-template main result \\
\path{traces_rerun_20260830_trained_80} & 80 & 9 trained & 5 & pass & trained-kernel active result \\
\path{traces_rerun_20260830_p2p_2000} & 2000 & 1 trained & 1 & pass & P2P-style generator \\
\path{traces_rerun_20260830_trained_avg1_120} & 120 & 9 trained & 1 & pass & harder stress test \\
\bottomrule
\end{tabular}
\end{table}

Table~\ref{tab:results} shows the reconstruction results. The all-zero pixel
baseline is around 0.88 because MNIST images have large black background. So
pixel accuracy alone is not enough. I also report foreground F1, MSSIM, and
recognition accuracy from the recovered image.

\begin{table}[H]
\centering
\caption{Reconstruction results from live CW305 traces.}
\label{tab:results}
\small
\begin{tabular}{>{\raggedright\arraybackslash}p{0.30\linewidth}cccccc}
\toprule
Method & Test N & Pixel acc. & FG F1 & MSSIM & Rec. bin & Rec. gray \\
\midrule
Active template, one-hot held-out & 40 & 0.9455 & 0.7495 & 0.726 & 0.9750 & -- \\
Active template, trained kernels & 40 & 0.9600 & 0.7924 & 0.769 & 0.9000 & -- \\
Hard active template, trained avg=1 & 20 & 0.9115 & 0.3770 & 0.452 & 0.4500 & -- \\
P2P clock-abs, gradient loss & 300 & 0.9503 & 0.7951 & 0.522 & 0.8933 & 0.9300 \\
P2P window-abs, gradient loss & 300 & 0.9486 & 0.7796 & 0.508 & 0.8333 & 0.8533 \\
P2P clock-abs, MSE loss & 300 & 0.9537 & 0.8068 & 0.561 & 0.9033 & 0.9133 \\
\bottomrule
\end{tabular}
\end{table}

\begin{figure}[H]
  \centering
  \includegraphics[width=\linewidth]{fig06_pilot_metrics.png}
  \caption{Fresh checked hardware results. The hard avg=1 run is included to
  show how much the result drops when repeat averaging is removed.}
  \label{fig:metrics}
\end{figure}

The one-hot active-template result is the cleanest proof that the trace carries
patch information. It reached 97.5\% recognition on 40 held-out images. The
trained-kernel active-template result is closer to a normal first convolution.
It has better pixel and structure scores, but lower recognition at 90.0\%.

\begin{figure}[H]
  \centering
  \includegraphics[width=0.62\linewidth]{fig15_rerun_onehot_heldout_examples.png}
  \caption{Active-template one-hot held-out examples.}
  \label{fig:onehot-examples}
\end{figure}

\begin{figure}[H]
  \centering
  \includegraphics[width=0.62\linewidth]{fig18_rerun_trained_heldout_examples.png}
  \caption{Active-template trained-kernel held-out examples.}
  \label{fig:trained-examples}
\end{figure}

The Power2Picture-style result is also important because it uses a learned
trace-to-image mapping. The best binary output came from the clock-abs MSE
model. The best gray recognition came from the clock-abs gradient-loss model.
I would keep both losses for a larger run.

\begin{figure}[H]
  \centering
  \includegraphics[width=0.82\linewidth]{fig17_rerun_p2p_examples.png}
  \caption{Power2Picture-style generator examples with clock-abs features and
  gradient-MSE loss.}
  \label{fig:p2p-grad}
\end{figure}

\begin{figure}[H]
  \centering
  \includegraphics[width=0.82\linewidth]{fig19_rerun_p2p_mse_examples.png}
  \caption{Power2Picture-style generator examples with clock-abs features and
  MSE loss.}
  \label{fig:p2p-mse}
\end{figure}

\section{LIMITATIONS AND CONCLUSION}

The main conclusion is that the experiment now has real hardware evidence. The
custom FPGA design was active, the functional output matched the CPU golden
model, and the traces were captured through ChipWhisperer. The reconstruction
methods used the side-channel traces at attack time. The truth images and
labels were used only for training/profile creation or for scoring after the
attack output existed.

The result is not a perfect reproduction of Power2Picture. The hardware is
different, and the measurement quality is lower. Also, the custom leakage bank
makes this a controlled victim design. This should be written clearly, because
the novelty here is not ``same hardware, same exact paper''. The novelty is
showing that a low-cost CW305/ChipWhisperer setup can still demonstrate
trace-to-image leakage from an FPGA neural-network style datapath, if the
gateware and capture method are adapted carefully.

The harder avg=1 test shows the current limit. With no repeat averaging, pixel
accuracy stayed above the all-zero baseline, but the digit recognizer dropped to
45\%. This means the single-shot trace still has some image information, but
not enough for a strong active-template reconstruction with the current search
and features.

The next work should:

\begin{itemize}
  \item scale the P2P hardware capture beyond 2,000 traces,
  \item add a build-ID register to make the hardware identity check cleaner,
  \item compare the leakage-bank design against a more natural accelerator,
  \item improve the active-template search speed for large held-out tests,
  \item test different clocks, gains, and averaging counts,
  \item keep the black/white hardware checklist in every run folder.
\end{itemize}

\begin{figure}[H]
  \centering
  \includegraphics[width=\linewidth]{fig08_next_hardware_checklist.png}
  \caption{Hardware checklist for future runs. This diagram is kept in
  black/white/gray because it is a flow chart.}
  \label{fig:checklist}
\end{figure}

\section{FILES BEHIND THIS REPORT}

The main files behind this report are:

\begin{itemize}
  \item \path{host/traces_rerun_20260830_onehot_80}
  \item \path{host/traces_rerun_20260830_trained_80}
  \item \path{host/traces_rerun_20260830_p2p_2000}
  \item \path{host/traces_rerun_20260830_trained_avg1_120}
  \item \path{report/professor_catchup_2026-08-30/validation_rerun_onehot_80.json}
  \item \path{report/professor_catchup_2026-08-30/validation_rerun_trained_80.json}
  \item \path{report/professor_catchup_2026-08-30/validation_rerun_p2p_2000.json}
  \item \path{report/professor_catchup_2026-08-30/validation_rerun_trained_avg1_120.json}
  \item \path{report/professor_catchup_2026-08-30/validation_extra_no_program_img1234_k4.json}
  \item \path{attack/rerun_20260830/score_onehot_80_eval40_heldout.json}
  \item \path{attack/rerun_20260830/score_trained_80_eval40_heldout.json}
  \item \path{attack/rerun_20260830/score_trained_avg1_120_eval20_heldout.json}
  \item \path{attack/rerun_20260830/p2p_2000_clockabs/summary.json}
  \item \path{attack/rerun_20260830/p2p_2000_windowabs/summary.json}
  \item \path{attack/rerun_20260830/p2p_2000_clockabs_mse/summary.json}
\end{itemize}

\end{document}
